% !TEX root = ../main.tex

% Give the second heading its own PDF destination, then resume the
% section and theorem counters where the first half of Talk 4 ended.
\edef\talkfourlastsection{\arabic{section}}
\edef\talkfourlasttheorem{\arabic{theorem}}
\addtocounter{chapter}{-1}
\let\talkfourHchapter\theHchapter
\renewcommand{\theHchapter}{\arabic{chapter}b}
\renewcommand{\thechapter}{\arabic{chapter}b}
\chapter{Derived manifolds}
\label[chapter]{chap:From_Smooth_Algebra_To_Derived_Manifolds}
\renewcommand{\thechapter}{\arabic{chapter}}
\setcounter{section}{\talkfourlastsection}
\setcounter{theorem}{\talkfourlasttheorem}

Consider again the origin $0\colon *\to\R$. We saw in
\Cref{chap:Ordinary_Finite_Limit_Completion_Manifolds} that its
self-intersection in the ordinary finite-limit completion is the point.
Could a different enlargement give a better answer? In any ordinary
category in which $*$ is terminal, the universal property forces
\[
  *\times_{\R}*\cong *.
\]
Indeed, for each object $T$ there is exactly one pair of maps $T\to *$,
and their composites with $0$ agree. Thus the point itself represents
the pullback. This argument works for a point in any ambient manifold.

The intersection of two points in a line has expected dimension
$0+0-1=-1$. Moving one copy of the origin slightly gives an empty
transverse intersection, so the intersection class vanishes. The ordinary
self-pullback, however, is a single point of dimension zero. A single
point represents a nonzero class in unoriented bordism, since the
boundary of a compact one-manifold has an even number of points.
Spivak uses this example to show that no ordinary category can satisfy
his general cup-product formula in cobordism
\cite[Definition~1.7 and Proposition~1.10]{Spivak_Derived_Smooth_Manifolds}.
We return to bordism in \Cref{thm:Spivak_Derived_Bordism_Comparison}.

The ordinary argument turns on what it means for the two composites
to agree. Suppose we could specify an identification between them
as part of a point of the pullback. The two maps to $*$ would still
be unique, but their common composite could have several
self-identifications. We first introduce the language for such a
pullback, then calculate the self-intersection of the origin and
formulate the universal property of derived manifolds.

\section{Pullbacks with higher identifications}
\label[section]{sec:Equality_As_Structure}

For sets, asking whether two elements are equal is a yes-or-no
question. For points of a space, one can instead ask for a path between
them, and then for a homotopy between two such paths. Animae give us
a language for these successive identifications.

\begin{definition}[Anima, informal description]
\label[definition]{def:Anima_Conceptual}
An \emph{anima} behaves like a set in which the identifications between
any two elements themselves form an anima. Identifications can be composed
and inverted, with the usual identities holding through coherent higher
identifications.
\end{definition}

We write $\An$ for the category of animae, understood in the higher
categorical sense described below. A set determines a discrete anima, a groupoid determines one whose
identifications are its isomorphisms, and a topological space $X$
determines its homotopy type $\Pi_\infty(X)$. In the last example,
paths give identifications and homotopies of paths give identifications
between them. For instance, the pullback of the base point of
$\Pi_\infty(S^1)$ along itself is its based loop anima. This is
equivalent to the discrete anima $\Z$ because each winding-number
component of the loop space is contractible.

In an $\infty$-category, the morphisms between two objects form an
anima. Composition is associative and unital with coherent higher
identifications. Ordinary categories are examples in which all the
hom animae are discrete. To form a pullback, we now apply the
universal property to these hom animae.
Let
\[
  X\xrightarrow{f}Z\xleftarrow{g}Y
\]
be a diagram in an $\infty$-category $C$. Its pullback is characterized by
natural equivalences of animae
\[
  \Hom_C(T,X\times_ZY)
  \simeq
  \Hom_C(T,X)
  \times_{\Hom_C(T,Z)}
  \Hom_C(T,Y)
\]
for all $T\in C$.

Informally, a generalized $T$-point consists of
\[
  a\colon T\to X,
  \qquad
  b\colon T\to Y,
  \qquad
  \alpha\colon f\circ a\cong g\circ b.
\]
The identification $\alpha$ is part of the datum. Equalities between choices of
$\alpha$, and all subsequent coherence, are contained in the pullback of hom
animae. In particular, a map from $T$ to the self-intersection of a
point $*\to X$ now amounts to a loop in $\Hom_C(T,X)$ based at the
composite $T\to *\to X$.

For a derived intersection the ambient $\infty$-category is the
opposite of a category of rings. We now describe its hom animae and
compute the self-intersection in that category.

\section{Animated \texorpdfstring{$\Cinfty$}{C-infinity}-rings and compact objects}
\label[section]{sec:Animated_Finite_Limit_Envelope}

The definition of a $\Cinfty$-ring object extends to any $\infty$-category
$C$ with finite products by taking $\operatorname{Alg}_{\Cinfty}(C)
=\Fun^\times(\mathrm{Euc},C)$. We apply it first to $C=\An$.

\begin{definition}[Animated $\Cinfty$-ring]
  \label[definition]{def:Animated_Cinfty_Ring}
  An \emph{animated $\Cinfty$-ring} is a finite-product-preserving functor
  \[
    A\colon\mathrm{Euc}\longrightarrow\An.
  \]
  We denote their $\infty$-category by
  \[
    \operatorname{Alg}_{\Cinfty}(\An)
    =
    \Fun^\times(\mathrm{Euc},\An).
  \]
\end{definition}

An ordinary $\Cinfty$-ring is the special case in which every value $A(\R^n)$
is a set, regarded as a $0$-truncated anima. Animated rings allow the smooth
operations themselves to interact with higher equality data.

Write $A\otimes^{\mathbb L}_{\infty,B}C$ for a pushout of animated
$\Cinfty$-rings.
For an animated ring $D$, the components of $D(\R)$ form its ordinary
ring $\pi_0D$. Based loops at zero up to homotopy give $\pi_1D$,
and the higher homotopy groups at zero give $\pi_iD$.

There is a useful calculation rule. For ordinary $\Cinfty$-rings
$A\leftarrow B\to C$, if one of the maps out of $B$ is surjective, then
\[
  \pi_i(A\otimes^{\mathbb L}_{\infty,B} C)
  \cong \operatorname{Tor}^{B}_i(A,C).
\]
On the right we take a free resolution of $A$ over the underlying
ordinary ring $B$ and tensor it over $B$ with $C$.

\begin{example}[The derived self-intersection of the origin]
\label[example]{ex:Animated_Self_Intersection_Origin}
For the diagram $*\to\R\leftarrow *$, the corresponding diagram of
rings is $\R\leftarrow B\to\R$, where $B=\Cinfty(\R)$ and both maps
are evaluation at zero. Multiplication by $x$ gives
a free resolution $0\to B\xrightarrow{x}B\to\R\to0$. Indeed, the ideal of
functions vanishing at zero is $(x)$ by Hadamard's lemma, and a smooth
function annihilated by $x$ vanishes on a dense open set. Tensoring
this resolution over $B$ with $\R$ gives the zero differential.
Thus the self-intersection has $\pi_0=\R$, $\pi_1=\R$, and all higher
groups vanish. For the derived zero locus of $t\mapsto t^2$,
tensoring the same resolution along $x\mapsto t^2$ gives multiplication
by $t^2$, which is injective. Here the infinitesimal information is
the nilpotent class of $t$ in $\pi_0=\Cinfty(\R)/(t^2)$.
\end{example}

The surjectivity hypothesis enters in passing to Tor. For a pushout
of animated rings along a map surjective on $\pi_0$, forgetting the
smooth operations gives a pushout of simplicial commutative
$\R$-algebras. This is the effective epimorphism case of
\citet[Theorem~3.32, Lemma~4.6 and Corollary~4.7]{Carchedi_Steffens_Universal_Derived_Manifolds}.
For ordinary inputs, the Tor spectral sequence for that pushout has
only its internal-degree-zero row, giving the formula above
\cite[Remark~4.2]{Carchedi_Steffens_Universal_Derived_Manifolds}.

\paragraph{Finite presentation for animated rings.}
For ordinary $\Cinfty$-rings, \Cref{prop:Filtered_Colimit_Criterion} expresses
finite presentation by preservation of filtered colimits. We use the same
criterion here, applied to the whole mapping anima.

\begin{definition}[Homotopical finite presentation]
  \label[definition]{def:Homotopically_Finitely_Presented_Cinfty_Ring}
  An animated $\Cinfty$-ring $A$ is \emph{homotopically finitely presented} if
  it is compact in $\operatorname{Alg}_{\Cinfty}(\An)$, meaning that the functor
  \[
    \Hom_{\operatorname{Alg}_{\Cinfty}(\An)}(A,-)
    \colon
    \operatorname{Alg}_{\Cinfty}(\An)
    \longrightarrow
    \An
  \]
  preserves filtered colimits.
\end{definition}

The free ring on $n$ generators is the discrete algebra $\Cinfty(\R^n)$.
Indeed its hom anima into $B$ is $B(\R^n)$ by Yoneda. Finite products of
animae commute with sifted colimits, so sifted colimits of algebras
are computed by evaluation. In particular these free rings are compact.

Write $\mathcal A=\operatorname{Alg}_{\Cinfty}(\An)$ and let
$\mathcal A_{\mathrm{fp}}$ denote its full subcategory of compact
objects. The category $\mathcal A$ is the sifted-colimit completion of
$\mathrm{Euc}\catop$. Its compact objects are precisely the retracts of
finite colimits of finitely generated free rings
\cite[Theorems~3.10--3.11 and Lemma~3.21]{Carchedi_Steffens_Universal_Derived_Manifolds}.
Indeed, finite colimits of compact objects are compact since filtered
colimits of animae commute with finite limits. Conversely, a presentation
of $A$ as a colimit of free rings expresses it as a filtered colimit of
colimits over finite simplicial subsets of the indexing simplicial set.
If $A$ is compact, its identity factors through one of these finite
colimits, exhibiting $A$ as a retract.

\begin{warning}
  \label[warning]{warn:Ordinary_Fp_Not_Necessarily_Animated_Fp}
  An ordinary finitely presented $\Cinfty$-ring need not remain compact when it
  is regarded as a discrete animated $\Cinfty$-ring. Thus the ordinary
  envelope is not simply a full subcategory of the animated envelope obtained
  by declaring every ordinary object discrete.
\end{warning}

The analogous warning for an arbitrary algebraic theory is recorded in
\citet[Warning~3.28]{Carchedi_Steffens_Universal_Derived_Manifolds}.
The truncation functor $\pi_0$ does send homotopically finitely presented
animated rings to finitely presented ordinary rings, but the reverse passage
can require retaining higher relation data.

\section{The universal property of derived manifolds}
\label[section]{sec:Derived_Manifolds_Universal_Pair}

The self-intersection has supplied the missing excess direction in
degree one. We now return to the question of
\Cref{chap:Ordinary_Finite_Limit_Completion_Manifolds} and allow the
target of a functor on manifolds to be an $\infty$-category.

\begin{definition}[Universal property of derived manifolds]
  \label[definition]{def:Universal_Infinity_Category_Derived_Manifolds}
  A pair $(D,i)$ has the universal property of \emph{derived manifolds}
  if $D$ is an idempotent-complete $\infty$-category with finite limits,
  $i\colon\mathrm{Mfd}\to D$ preserves transverse pullbacks and the
  terminal object, and restriction along $i$ induces an equivalence
  \[
    \Fun^{\mathrm{lex}}(D,C)\simeq\Fun^{\pitchfork}(\mathrm{Mfd},C)
  \]
  for every idempotent-complete $\infty$-category $C$ with finite
  limits. The left side consists of finite-limit-preserving functors,
  the right side of functors preserving transverse pullbacks and the
  terminal object. For such a pair we write $\mathrm{DMfd}$ for $D$.
\end{definition}

Here idempotent completeness means that coherent idempotents split,
as explained in \Cref{sec:From_Smooth_Algebra_Summary_Supplementary_Materials}.
The following theorem constructs a pair with this universal property.

\begin{theorem}[Derived manifolds from animated $\Cinfty$-rings]
  \label[theorem]{thm:Universal_Property_Derived_Manifolds_Cinfty_Rings}
  The functor
  $i\colon\mathrm{Mfd}\to(\mathcal A_{\mathrm{fp}})\catop$ taking $M$
  to its discrete ring $\Cinfty(M)$ is fully faithful, and the pair
  $((\mathcal A_{\mathrm{fp}})\catop,i)$ has the universal property
  of derived manifolds.
\end{theorem}

\begin{corollary}
  \label[corollary]{cor:Derived_Manifolds_Animated_Cinfty_Rings}
  There is a canonical equivalence under $\mathrm{Mfd}$
  \[
    \mathrm{DMfd}\simeq\operatorname{Alg}_{\Cinfty}(\An)_{\mathrm{fp}}\catop .
  \]
\end{corollary}

The corollary follows from uniqueness of the universal pair. For the
theorem we must answer two questions. Algebraically, why do compact
animated rings give the finite limits generated by Euclidean spaces?
Geometrically, why does a transverse pullback still give the correct
pushout after animation? Applying $\pi_0$ to an animated pushout gives
the ordinary pushout, since $\pi_0$ is left adjoint to the inclusion of
discrete rings. Thus \Cref{thm:Transverse_Pullbacks_Pushouts} already
identifies $\pi_0$. We will calculate the higher homotopy groups for a regular
zero locus, then reduce an arbitrary transverse pullback to that case.
The universal property follows by extending functors from Euclidean
spaces to manifolds
\cite[Definition~2.1, Theorem~5.3 and Corollary~5.4]{Carchedi_Steffens_Universal_Derived_Manifolds}.

\section{The animated finite-limit envelope}
\label[section]{sec:Animated_Envelope_Proof}

Let $j^{\mathrm{fp}}\colon\mathrm{Euc}\to(\mathcal A_{\mathrm{fp}})\catop$
be the inclusion of the free rings, with variance reversed. We will
extend a functor on Euclidean spaces by first allowing it to take
values in presheaves and then using compactness to show that these
presheaves are representable.

\begin{theorem}[Animated finite-limit envelope]
  \label[theorem]{thm:Animated_Cinfty_Finite_Limit_Envelope}
  For every idempotent-complete $\infty$-category $C$ with finite
  limits, restriction along $j^{\mathrm{fp}}$ gives an equivalence
  \[
    \Fun^{\mathrm{lex}}\bigl((\mathcal A_{\mathrm{fp}})\catop,C\bigr)
    \simeq\Fun^\times(\mathrm{Euc},C).
  \]
\end{theorem}

\begin{proof}
  Let $F\colon\mathrm{Euc}\to C$ preserve finite products. For $c\in C$
  the functor $B_c(T)=\Hom_C(c,F(T))$ is an animated $\Cinfty$-ring.
  To each $A\in\mathcal A$ assign the presheaf on $C$ given by
  \[
    \widetilde F(A)(c)=\Hom_{\mathcal A}(A,B_c).
  \]
  This assignment takes colimits in $A$ to limits of presheaves. On
  $A=\Cinfty(T)$ its value at $c$ is $B_c(T)=\Hom_C(c,F(T))$ by Yoneda,
  so it extends the composite of $F$ with the Yoneda embedding $y$.
  More precisely, consider the category of finitely generated free rings
  over $A$. It has finite coproducts, formed by taking the coproduct of
  the free rings with its induced map to $A$, so it is sifted.
  By Yoneda density, the colimit of this diagram of representables is
  the presheaf $A$. Sifted colimits of animated $\Cinfty$-rings are
  computed pointwise, so this is also a colimit in $\mathcal A$.
  Mapping it into $B_c$ gives the limit over the opposite indexing
  category, exactly the pointwise right Kan extension formula for $yF$.

  When $A$ is compact, it is a retract of a finite colimit of free rings.
  Hence $\widetilde F(A)$ is a retract of a finite limit of representables.
  The image of $y$ is closed under both operations since $C$ has finite
  limits and is idempotent-complete. The formula therefore defines a
  functor $\widehat F\colon(\mathcal A_{\mathrm{fp}})\catop\to C$,
  left exact because $y$ detects limits. It is the right Kan extension
  of $F$ along $j^{\mathrm{fp}}$ since the pointwise limits just computed
  in presheaves are represented in $C$.

  For a left-exact $G$ on $(\mathcal A_{\mathrm{fp}})\catop$, the unit
  $G\to\widehat{G j^{\mathrm{fp}}}$ is an equivalence on free rings by
  full faithfulness of $j^{\mathrm{fp}}$. Both functors preserve finite
  limits and retracts, so it is an equivalence everywhere. The counit
  is evaluation on free rings, already identified by Yoneda. These
  natural maps give the asserted equivalence of functor categories.
\end{proof}

This is \citet[Theorem~3.22]{Carchedi_Steffens_Universal_Derived_Manifolds}.
To use it for manifolds we must show that their function rings are
compact and identify the effect of the functor on transverse pullbacks.

\section{Regular zero loci and compactness of function rings}
\label[section]{sec:Animated_Regular_Zero_Loci}

Let $g=(g_1,\ldots,g_q)\colon M\to\R^q$ be smooth and write $Z=g^{-1}(0)$.
When $0$ is a regular value, \Cref{cor:Function_Ring_Regular_Zero_Locus}
identifies $\Cinfty(Z)$ with the quotient by $(g_1,\ldots,g_q)$.
We now compute the positive homotopy groups of the same pushout in
animated rings.

\begin{proposition}[Derived transverse zero locus]
  \label[proposition]{prop:Animated_Regular_Zero_Locus}
  If $0$ is a regular value of $g$, restriction induces an equivalence
  of animated $\Cinfty$-rings
  \[
    \Cinfty(M)\otimes^{\mathbb L}_{\infty,\Cinfty(\R^q)}\R
    \longrightarrow \Cinfty(Z).
  \]
\end{proposition}

\begin{proof}
  Write $B=\Cinfty(\R^q)$. The evaluation $B\to\R$ is surjective,
  so we may compute the homotopy groups by Tor. The coordinates
  $x_1,\ldots,x_q$ form a regular sequence in $B$. Indeed, quotienting
  by the first $j$ coordinates gives the functions on the corresponding
  coordinate subspace by Hadamard's lemma, and multiplication by the
  next coordinate is injective there. At each step, adjoining the next
  two-term Koszul factor takes the mapping cone of multiplication by
  that coordinate. The long exact homology sequence gives only the
  next quotient in degree zero. Thus the Koszul complex on all $q$
  coordinates is a finite free resolution of $\R$ over $B$.
  Tensoring it with $\Cinfty(M)$ replaces $x_j$ by $g_j$. Its degree
  $p$ term is $\Cinfty(M)\otimes_{\R}\Lambda^p((\R^q)^\vee)$, with
  differential contraction by $g$. Write $e_j$ for the $j$th coordinate
  covector, so that $d(e_j)=g_j$.

  To compute its homology, take the same complex of sheaves on $M$.
  Near a point of $Z$, the submersion theorem gives coordinates in
  which $g_1,\ldots,g_q$ are the first $q$ coordinates. The coordinate
  calculation identifies the homology there with smooth functions on
  $Z$ in degree zero. Off $Z$, some $g_j$ is locally invertible and
  exterior multiplication by $e_j/g_j$ contracts the complex.
  These local identifications use restriction to $Z$ as augmentation.
  Global sections are exact on sheaves of smooth-function modules on
  $M$ since local lifts of a section can be multiplied by a subordinate
  partition of unity and summed to give a global lift. The global
  complex therefore has homology $\Cinfty(Z)$ in degree zero and zero
  in positive degrees. This proves that the comparison is an
  equivalence.
\end{proof}

This module calculation is \citet[Lemma~4.5]{Carchedi_Steffens_Universal_Derived_Manifolds}.
In \Cref{chap:Derived_Zero_Loci_And_Tangent_Complexes} we will construct
a differential graded $\Cinfty$-algebra presenting the pushout itself.

\begin{corollary}[Compactness of manifold function rings]
  \label[corollary]{cor:Manifold_Function_Ring_Animated_Compact}
  For every manifold $M$, its discrete ring $\Cinfty(M)$ is compact in
  $\mathcal A$.
\end{corollary}

\begin{proof}
  For an open $U\subseteq\R^d$, choose a smooth function $a$ with
  $U=\{a\ne0\}$ by \Cref{lem:Open_Subset_Cozero_Locus}. Projection
  identifies the zero locus of $h(x,y)=a(x)y-1$ with $U$. On that
  locus $\partial_yh=a(x)\ne0$, so
  \Cref{prop:Animated_Regular_Zero_Locus} presents $\Cinfty(U)$ as
  a pushout of three free rings, $\Cinfty(\R^{d+1})$, $\Cinfty(\R)$
  and $\R$. Thus it is compact. For general $M$, take a proper embedding
  in Euclidean space. The retraction of an open tubular neighbourhood
  onto $M$ makes $\Cinfty(M)$ a retract of a compact object.
\end{proof}

We use this same geometric presentation below
\cite[Proposition~4.9 and Corollary~4.10]{Carchedi_Steffens_Universal_Derived_Manifolds}.
In particular the functor $i$ in the main theorem is now defined on all
manifolds. It is fully faithful since ordinary rings embed fully faithfully
as discrete animated rings and \Cref{cor:Recovering_Manifolds} identifies
maps between manifold function rings.

\section{Derived transversality}
\label[section]{sec:Derived_Transversality}

We next pass from a regular zero locus to a transverse intersection
$P=M\times_LN$. The product case is needed first, because the difference
of two maps to $\R^d$ is defined on $M\times N$.

\begin{lemma}
  \label[lemma]{lem:Animated_Manifold_Products}
  The functor $i$ preserves finite products.
\end{lemma}

\begin{proof}
  The point goes to the initial ring $\R$ in $\mathcal A$. For
  Euclidean spaces this follows because the coproduct
  of the free rings on $m$ and $n$ generators is free on $m+n$.
  For Euclidean opens $U$ and $V$, use the scalar equations
  $h_U=0$ and $h_V=0$ from the preceding proof. The two equations on
  the product of the ambient Euclidean spaces have independent
  differentials along $U\times V$, since they use disjoint variables.
  Taking the coproduct of the two pushout presentations gives the
  pushout for $(h_U,h_V)$, since the coproducts of the three respective
  pairs of free rings are the free rings for the product ambient
  space, $\R^2$, and the point. By
  \Cref{prop:Animated_Regular_Zero_Locus} this presents
  $\Cinfty(U\times V)$, with structure maps given by the projections.
  Finally, for $M$ and $N$ choose tubular neighbourhoods $U$ and $V$
  retracting onto them. By naturality, the product comparison for
  $M,N$ is a retract of the comparison for $U,V$, hence an equivalence.
\end{proof}

\begin{theorem}[Preservation of transverse pullbacks]
  \label[theorem]{thm:Animated_Transverse_Pullbacks}
  If $f\colon M\to L$ and $g\colon N\to L$ are transverse, then
  restriction along the two projections induces an equivalence
  \[
    \Cinfty(M)\otimes^{\mathbb L}_{\infty,\Cinfty(L)}\Cinfty(N)
    \longrightarrow \Cinfty(M\times_L N).
  \]
\end{theorem}

\begin{proof}
  We prove that $i$ carries the transverse square to a pullback, in
  three steps according to the target $L$.

  \emph{Euclidean target.}
  Suppose $L=\R^d$. The difference $h(m,n)=f(m)-g(n)$ has zero locus
  $P$, and $dh$ is surjective there by transversality. By
  \Cref{prop:Animated_Regular_Zero_Locus}, $i(P)$ is the zero fibre of
  $i(h)$. To identify this fibre with the required pullback, observe
  that on $\R^d\times\R^d$ the automorphism $(v,w)\mapsto(v-w,w)$
  carries the diagonal to $\{0\}\times\R^d$. By
  \Cref{lem:Animated_Manifold_Products} the same automorphism identifies
  the pullback of the diagonal along $(i(f),i(g))$ with the zero fibre
  of $i(h)$ on $i(M)\times i(N)$.

  \emph{Open Euclidean target.}
  Let $U\subseteq\R^d$ be open. Applying the first step to the open
  inclusion twice gives
  $i(U)\simeq i(U)\times_{i(\R^d)}i(U)$ via the diagonal. Thus
  $i(U)\to i(\R^d)$ is a monomorphism. For maps landing in $i(U)$
  we may consequently compute their pullback over $i(\R^d)$.
  Transversality of maps into $U$ is unchanged by the open inclusion,
  so the first step proves the assertion over $U$.

  \emph{Manifold target.}
  Embed $L$ in $\R^d$ and let $e\colon L\to U$ and $r\colon U\to L$
  be the zero section and projection of an open tubular neighbourhood.
  Since $r$ is a submersion, the ordinary manifold $N'=N\times_L U$
  exists. Write $g'\colon N'\to U$ for its second projection. Consider
  the squares
  \[
  \begin{tikzcd}
    P \rar \dar & N \dar["g"] \\
    M \rar["f"'] & L
  \end{tikzcd}
  \qquad\qquad
  \begin{tikzcd}
    P \rar \dar & N' \dar["g'"] \\
    M \rar["ef"'] & U .
  \end{tikzcd}
  \]
  The square on the right is a transverse pullback. Indeed, over
  $e(l)$ the differential of $g'$ supplies all directions in
  $\ker d r_{e(l)}$, and its projection to $T_lL$ has image
  $\operatorname{im}dg$. Together with $d(ef)$ this spans $T_{e(l)}U$
  by the original transversality. A point in the new pullback has
  the form $(m,(n,e(f(m))))$ with $f(m)=g(n)$, and this identification
  and its inverse are smooth.

  The left square is a retract of the right square as a diagram of
  manifolds. On the bottom right use $e$ and $r$, on the top right use
  $n\mapsto(n,e(g(n)))$ and projection to $N$, and on $M$ and $P$ use
  the identities. Applying $i$ preserves this retraction. The square
  over $i(U)$ is cartesian by the second step, so its retract is
  cartesian as well, since its comparison to the pullback is a retract of
  an equivalence.
\end{proof}

The product lemma is \citet[Lemma~4.11]{Carchedi_Steffens_Universal_Derived_Manifolds},
and the theorem is their Theorem~4.52. Their proof uses completeness of
the higher homotopy groups as modules over $\pi_0$ to pass from local
calculations to global sections (Theorem~4.44). The proof here uses
the tubular retraction of squares for this passage. The manifold $N'$
in the proof is formed in $\mathrm{Mfd}$ before applying $i$.

\section{Extension from Euclidean spaces to manifolds}
\label[section]{sec:Universal_Property_Derived_Manifolds}

Let $q\colon\mathrm{Euc}\to\mathrm{Mfd}$ be the inclusion, so that
$i q=j^{\mathrm{fp}}$. The algebraic envelope gives an extension
$\widehat F$ of any product-preserving $F$ on $\mathrm{Euc}$. Derived
transversality now lets us restrict that extension to manifolds.

\begin{proposition}[Geometric extension]
  \label[proposition]{prop:Geometric_Extension_From_Euclidean_Spaces}
  For every idempotent-complete $\infty$-category $C$ with finite limits,
  restriction along $q$ induces an equivalence
  \[
    \Fun^{\pitchfork}(\mathrm{Mfd},C)\simeq\Fun^\times(\mathrm{Euc},C).
  \]
  Its inverse sends $F$ to $\widehat F\circ i$.
\end{proposition}

\begin{proof}
  Restriction is well defined since products of manifolds are transverse
  pullbacks over the point. The composite $\widehat F\circ i$ preserves
  transverse pullbacks and the terminal object by left exactness of
  $\widehat F$, \Cref{thm:Animated_Transverse_Pullbacks} and
  \Cref{lem:Animated_Manifold_Products}. It is the pointwise right Kan
  extension of $F$ along $q$. Indeed, at $M$ the indexing category for
  this extension consists of maps $M\to q(T)$, for $T$ Euclidean.
  Full faithfulness of $i$ identifies it with the indexing category
  of maps $i(M)\to j^{\mathrm{fp}}(T)$ computing $\widehat F(i(M))$.

  The counit on Euclidean spaces is an equivalence since $q$ is fully
  faithful. For $G\in\Fun^{\pitchfork}(\mathrm{Mfd},C)$ consider the
  unit $G\to\operatorname{Ran}_q(Gq)$. It is an equivalence on
  Euclidean spaces. Each Euclidean open is a transverse pullback of
  Euclidean spaces by its presentation $a(x)y=1$, and both functors
  preserve that pullback. Thus the unit is an equivalence on Euclidean
  opens, and then on every manifold by the tubular retraction. This
  proves the equivalence on functor categories
  \cite[Lemma~5.1 and Corollary~5.2]{Carchedi_Steffens_Universal_Derived_Manifolds}.
\end{proof}

\begin{proof}[Proof of \Cref{thm:Universal_Property_Derived_Manifolds_Cinfty_Rings}]
  The compact algebras are closed under finite colimits and retracts,
  so their opposite has finite limits and is idempotent-complete.
  The functor $i$ is fully faithful, preserves the terminal object,
  and preserves transverse pullbacks, as proved above. For every
  idempotent-complete $C$ with finite limits, restriction along $i$ fits
  into the triangle
  \[
  \begin{tikzcd}[column sep=small]
    \Fun^{\mathrm{lex}}\bigl((\mathcal A_{\mathrm{fp}})\catop,C\bigr)
      \ar[rr,"i^*"] \ar[dr,"(j^{\mathrm{fp}})^*"']
    && \Fun^{\pitchfork}(\mathrm{Mfd},C) \ar[dl,"q^*"] \\
    & \Fun^\times(\mathrm{Euc},C).
  \end{tikzcd}
  \]
  The triangle commutes since $i q=j^{\mathrm{fp}}$. Its diagonal maps
  are equivalences by \Cref{thm:Animated_Cinfty_Finite_Limit_Envelope,prop:Geometric_Extension_From_Euclidean_Spaces}.
  Hence $i^*$ is an equivalence. This constructs the universal pair
  of \Cref{def:Universal_Infinity_Category_Derived_Manifolds}.
\end{proof}

In particular, for an arbitrary smooth map $f\colon M\to\R^q$, the
ring of the derived zero locus is the animated pushout
$\Cinfty(M)\otimes^{\mathbb L}_{\infty,\Cinfty(\R^q)}\R$.
The next chapter gives a multiplicative model for this ring and uses
it to compare singular intersections and their families.

\section{Further details on retracts and coherence}
\label[section]{sec:From_Smooth_Algebra_Summary_Supplementary_Materials}

The universal property of derived manifolds requires idempotent-complete targets.
We explain more precisely what this means in an $\infty$-category and
why finite limits alone do not supply it. The exercises revisit this
point and the self-pullback example.

Informally, an idempotent looks like a morphism $e\colon X\to X$ together with
a homotopy from $e\circ e$ to $e$. Formally, an idempotent in an
$\infty$-category is a coherent idempotent diagram: the displayed homotopy is
part of the data, together with its higher compatibilities. Splitting it means
finding an object $Y$ and maps
\[
  Y\xrightarrow{i}X\xrightarrow{r}Y
\]
together with coherent identifications of $r\circ i$ with
$\operatorname{id}_Y$ and of $i\circ r$ with the given idempotent. In an
ordinary category with finite limits the equalizer construction in
\Cref{sec:Ordinary_Finite_Limit_Envelope} supplies a splitting. For $\infty$-categories
this is not forced by finite limits alone. The compact algebraic objects are
closed under retracts, so their opposite category is naturally
idempotent-complete. The universal property of derived manifolds must record
this closure explicitly.

\begin{exercise}
Describe the equalizer splitting for an idempotent on an object of
$\operatorname{Aff}_{\Cinfty}^{\mathrm{fp}}$. Compare it with the
explicit finite presentation for a retract ring
in \Cref{chap:Manifolds_As_Cinfty_Rings}.
\end{exercise}

\begin{exercise}
Use the universal property of a pullback to identify
$*\times_{BG}*$ with a discrete group $G$, where $BG$ is its
classifying anima and both maps select the base point.
Compare this with the self-pullback over a set.
\end{exercise}

The primary source is
\cite[Definition~2.1, Theorem~3.22, Remark~3.26,
and Section~5]{Carchedi_Steffens_Universal_Derived_Manifolds}.
It treats both universal properties, their geometric comparison, and
the relation with Spivak's construction.

\let\theHchapter\talkfourHchapter
